\section{拉格朗日子空间与全实子空间}
\label{sec:lagrangian-versus-totally-real-subspaces}

\stfhIndex{I171}{Quanshizikongjian}{全实子空间（totally real subspace）}
为研究\termRiemannRochFormula{}的有边界版本，我们考虑 $\mathbb C^n$ 中（极大的）\emph{\termTotallyRealSubspace{}}。若实子空间 $V\subset\mathbb C^n$ 满足
\[
  V\cap iV=\{0\},
  \qquad
  \dim_{\mathbb R}V=n,
\]
则称 $V$ 是全实的。我们以 $\mathcal R(n)$ 表示所有全实子空间组成的集合。

任一全实子空间 $V$ 都可写成
\[
  V=A\cdot\mathbb R^n
\]
的形式，其中 $A\in GL(n,\mathbb C)$；而
$A_1\cdot\mathbb R^n=A_2\cdot\mathbb R^n$ 当且仅当
\begin{equation}
  A_2^{-1}A_1\in GL(n,\mathbb R).
\label{eq:totally-real-coset-condition}
\end{equation}
因此，全实子空间的集合 $\mathcal R(n)$ 是\termHomogeneousSpace{}：
\[
  \mathcal R(n)=GL(n,\mathbb C)/GL(n,\mathbb R).
\]
下面的引理是研究指标问题时一个有用的事实。其证明见 (Oh95b) 中的命题 4.4。

\begin{lemma}
\label{lem:totally-real-subspace-matrix-model}
考虑子集
\[
  \widetilde{\mathcal R}(n)
  :=\{D\in GL(n,\mathbb C)\mid D\overline D=I_n\},
\]
其中 $I_n$ 是单位矩阵。则映射
\[
  B:\mathcal R(n)\cong GL(n,\mathbb C)/GL(n,\mathbb R)
  \longrightarrow\widetilde{\mathcal R}(n);
  \qquad
  A\cdot\mathbb R^n\longmapsto A\overline A^{-1}
\]
是微分同胚。
\end{lemma}

\begin{proof}
考虑定义在 $GL(n,\mathbb C)$ 上的映射
\begin{equation}
  A\longmapsto A\overline A^{-1}.
\label{eq:matrix-conjugation-map}
\end{equation}
注意，任意形如 $A\overline A^{-1}$ 的矩阵 $D$ 都满足
$D\overline D=\operatorname{Id}_n$，所以这个映射的像落在
$\widetilde{\mathcal R}(n)$ 中。

方程~\eqref{eq:totally-real-coset-condition}等价于
\[
  A_1\overline A_1^{-1}=A_2\overline A_2^{-1}.
\]
这说明映射 $B$ 是良定义的。映射~\eqref{eq:matrix-conjugation-map}给出一一对应的证明留作练习。
\end{proof}

\begin{exercise}
\label{ex:matrix-conjugation-bijection}
证明映射~\eqref{eq:matrix-conjugation-map}给出一一对应。
\end{exercise}

下面考虑标准辛向量空间 $(\mathbb R^{2n},\omega_0)$，其中
\[
  \omega_0=\sum_{i=1}^n dx_i\wedge dy_i,
\]
这里使用典范坐标 $(x_1,\ldots,x_n,y_1,\ldots,y_n)$。以
\[
  \Lambda(n):=\mathcal L(\mathbb R^{2n},\omega_0)
\]
表示 $(\mathbb R^{2n},\omega_0)$ 的所有\termLagrangianSubspace{}组成的集合。若在 $\mathbb R^{2n}$ 上配备标准\termComplexMultiplication{}，并借映射
\[
  (x_i,y_i)\longmapsto z_i=x_i+\sqrt{-1}y_i
\]
将它与 $\mathbb C^n$ 等同，则每个拉格朗日子空间
$V\subset\mathbb C^n$ 都可以写成
\[
  V=A\cdot\mathbb R^n,
\]
其中复矩阵 $A\in U(n)$。这说明 $\Lambda(n)$ 是\termHomogeneousSpace{}：
\[
  \Lambda(n)\cong U(n)/O(n).
\]
特别地，这说明自然包含 $\Lambda(n)\hookrightarrow\mathcal R(n)$ 是\termHomotopyEquivalence{}。

(Ar67) 证明了
\[
  H^1(\Lambda(n),\mathbb Z)\cong\mathbb Z,
\]
而且 $\Lambda(n)$ 带有著名的\termCharacteristicClass{}
$\mu\in H^1(\Lambda(n),\mathbb Z)$，即\emph{\termMaslovClass{}} (Ar67)。它给每条给定的环路
$\gamma:S^1\to\Lambda(n)$ 指派如下整数\footnote{译者注：原书把 $\mu(\gamma)$ 写成 $\frac{1}{2\pi}\det^2(A)$；但右端按通常记号不是环路的整数值，也没有写出度数或辐角变化。译文照录原式，不自行重构。}：
\[
  \mu(\gamma)=\frac{1}{2\pi}\det^2(A).
\]
关于这个类的无坐标描述，参见第 3.2 节。此外，两条环路
$\gamma_1,\gamma_2$ 同伦，当且仅当
$\mu(\gamma_1)=\mu(\gamma_2)$。而且，只要双线性形式
$\omega_0(\,\cdot\,,J_0\,\cdot\,)$ 保持为\emph{正}，$\mu(\gamma)$ 的值就不依赖于 $\mathbb R^{2n}$ 与 $\mathbb C^n$ 的等同；这里 $J_0$ 是与 $\mathbb C^n$ 上复乘法相伴的殆复结构。

注意，当 $A$ 是\termUnitaryMatrix{}时，
$\overline A^{-1}=A^{\mathsf T}$，其中 $A^{\mathsf T}$ 是 $A$ 的转置。因此，由\cref{lem:totally-real-subspace-matrix-model}得到下面的推论。

\begin{corollary}
\label{cor:symmetric-unitary-model-lagrangian-grassmannian}
以 $\widetilde\Lambda(n)\subset\widetilde{\mathcal R}(n)$ 表示 $B$ 限制到
$\Lambda(n)\subset\mathcal R(n)$ 后的像，即
\begin{equation}
  \widetilde\Lambda(n)
  :=\{D\in U(n)\mid D=D^{\mathsf T}\},
\label{eq:symmetric-unitary-matrices}
\end{equation}
也就是\termSymmetricUnitaryMatrix{}组成的集合。则 $B$ 在 $\Lambda(n)$ 上的限制给出一个微分同胚。
\end{corollary}

下面把\termMaslovIndex{}推广到全实子空间的环路上；限制到
$\Lambda(n)$ 中的环路时，它就是马斯洛夫指标。

\begin{definition}[广义马斯洛夫指标]
\label{def:generalized-maslov-index}
设 $\gamma:S^1\to\mathcal R(n)$ 是一条环路。\termGeneralizedMaslovIndex{}
$\mu(\gamma)$ 定义为映射
\[
  \det\circ B\circ\gamma:S^1\longrightarrow\mathbb C\setminus\{0\}
\]
的\termWindingNumber{}。
\end{definition}
